

\section{Related Work}
\label{sec:related}

\noindent
Treemaps are a common technique for space-restricted visualization of tree-structured data. The original algorithm presented by Shneiderman recursively splits parent nodes alternating in horizontal and vertical direction by their child nodes~\cite{Shneiderman1992}, resulting in a 2D layout.
By using specific shading~\cite{VanWijk1999}, margins or insets~\cite{Lue2008}, the treemap can depict topology as well.
Often used visual variables include size and color~\cite{Fekete2002} as well as texture~\cite{Holten2005}.

\paragraph{Treemap Layouts}

Besides treemap presentations, a basic layout algorithm has to be chosen according to the particular use cases.
For example, the choice of the layout algorithm influences the creation and preservation of a mental map for the user of the visualization~\cite{Misue1995}.
One major disadvantage of the Slice'n Dice algorithm is the creation of unfavorable aspect ratios of the node representations. Bruls et~al. as well as Bederson and Shneiderman counterbalance this problem using Squarified~\cite{Bruls1999} and Strip Treemap~\cite{Bederson2002} layouts.
Latest research focuses on the use of space-filling curves and further optimizes the trade-off between readability and stability~\cite{Tu2007,Tak2013}.
Since most of the treemap layout algorithms are designed for 2D presentation, they do not take possible occlusion within 2.5D treemaps into account.

\paragraph{2.5D Treemaps}

The first extension of 2D treemap algorithms to the third dimension is presented in the StepTree visualization system by Bladh et~al.~\cite{Bladh2004}. It extends the initial approach by stacking the graphical representations of subdirectories.
In a comparative study between the StepTree and its corresponding 2D treemap, users performed significantly better in tasks of interpreting the hierarchical structure while preserving performance in other interpretive and navigational tasks.
In addition to this approach, there exist a variety of modifications and extensions, e.g., 3D Polar Treemaps, Treecubes, or Collector's Box Treemaps~\cite{Schulz2011}.
In contrast to general 3D treemaps, most current 3D tree visualization techniques combine a 2D layout and extrude the graphical elements forming the third dimension.
Due to the restricted use of the third dimension, we refer to the term 2.5D treemaps.
An example of how to adaptively extend the nodes' shape in 2D treemaps by the third dimension was provided by Chaudhuri and Shen~\cite{Chaudhuri2009}.
Wettel and Lanza use 2.5D treemaps to create a metaphorical view on software system artifact hierarchies with CodeCities~\cite{Wettel2007} by mapping additional information to the height of leaf nodes.
Bohnet and D{\"o}llner mentioned the benefits of using realistic rendering effects (e.g., shadows) in 2.5D treemaps to support users in distinguishing individual elements~\cite{Bohnet2011}.

The increasing number of items in large hierarchies leads to perceptional questions about treemap visualizations, e.g., how much effect does the width of a node's border have on the ability to extract the hierarchical structure of the underlying data~\cite{Kong2010}.
With an increasing number of items, the need for focus+context as well as overview+detail techniques emerges~\cite{Hauser2006}.
Even though the visualization of a wide range of large hierarchical datasets (e.g., software systems, economy data) is frequently studied, the interaction techniques used in this field remain straightforward.
Apparent interaction techniques in 2.5D treemaps are the selection of inner nodes and leaves as well as the coloring/highlighting of nodes~\cite{Wettel2007}.
Sud et~al. present interactive rendering, navigation, and animation of dynamically changing 2D Voronoi Treemaps by using a GPGPU-based rendering technique~\cite{Sud2010}.

\paragraph{Mixed-perspective Visualization}

The choice for the type of projection for virtual 3D environments is often either an orthographic or perspective projection~\cite{Piringer2004}.
However, depending on the composition and the characteristics of the environment, mixed projections can provide effective communication of the content.
For virtual terrain and city model visualization, multi-perspective views with bent projections provide a near-field depiction of high detail and a context area of less detail~\cite{Pasewaldt2011,Roehlig2016}.
Before hardware-accelerated rendering techniques for mixed projections were developed~\cite{Lorenz2008}, these approaches were already applied by artists such as H.~C.~Berann~\cite{Berann1999}.
